% !TEX root = ../algorithms.tex

\section{Модели вычислений: BSP и модель внешней памяти}

\begin{Definition}{Машинное слово}{}
	Машинное слово (в моделях, подобных RAM) --- это целочисленный тип данных, способный принимать все значения, необходимые для работы алгоритма.
	Машинное слово состоит из некоторого числа бит.
\end{Definition}

\subsection{Модель блочно-синхронного параллелизма (\emph{Bulk Synchronous Parallel, BSP})}

Модель \textbf{BSP} задаётся параметрами \(p\), \(g\) и \(L\), где:
\begin{itemize}
	\item \(p\) --- число компьютеров (процессоров);
	\item \(g\) --- время передачи или получения одного машинного слова;
	\item \(L\) --- время синхронизации.
\end{itemize}

Обычно предполагается, что \(g \gg 1,\quad L \gg 1.\)

\begin{enumerate}
	\item Модель состоит из \(p\) компьютеров (или процессоров), пронумерованных числами от \(1\) до \(p\).
	Каждый из них представляет собой полноценную RAM-модель со своей локальной памятью.
	
	\item Компьютеры объединены сетью, и каждый из них может посылать сообщения размера \(m\) машинных слов другому компьютеру за \(gm\) единиц времени.
	
	\item После этапа локальных вычислений все компьютеры
	``синхронизируются'', ожидая завершения работы всех остальных
	компьютеров. На этом этапе доставляются все отправленные сообщения.
	Каждый компьютер, получающий \(m\) машинных слов, тратит \(gm + L\)
	единиц времени: \(gm\) на получение сообщения и \(L\) на синхронизацию.
\end{enumerate}

\subsection{Супершаги}

Алгоритм в модели BSP выполняется по \emph{супершагам}. Каждый супершаг состоит из трёх последовательных этапов:
\begin{enumerate}
	\item локальные вычисления на каждом компьютере;
	\item обмен сообщениями между компьютерами;
	\item синхронизация, на которой все компьютеры дожидаются друг друга.
\end{enumerate}

Обозначим через \(t_k^{(s)}\) время локальных вычислений, выполненных \(k\)-м компьютером на \(s\)-м супершаге.
Через \(c_k^{(s)}\) обозначим количество машинных слов, которые \(k\)-й компьютер отправил или получил на этом супершаге.

Тогда время выполнения \(s\)-го супершага определяется самым медленным вычислением, самой большой коммуникационной нагрузкой и временем синхронизации. Поэтому оно равно
\[
\max_{k=1,\dots,p} t_k^{(s)}
\;+\;
\max_{k=1,\dots,p} \bigl(g\,c_k^{(s)}\bigr)
\;+\;
L.
\]
\subsection{Стоимость BSP-алгоритма}

Пусть алгоритм работает \(\hat S\) супершагов. Тогда его время работы в модели BSP удобно измерять тремя компонентами:
\begin{enumerate}
	\item \textbf{вычисления}
	\[
	\sum_{s=1}^{\hat S} \max_{k=1,\dots,p} t_k^{(s)};
	\]
	
	\item \textbf{коммуникация}
	\[
	\sum_{s=1}^{\hat S} \max_{k=1,\dots,p} c_k^{(s)};
	\]
	
	\item \textbf{синхронизация}
	\[
	\hat S\,L.
	\]
\end{enumerate}

Во второй компоненте формально должен стоять множитель \(g\), то есть
\[
\sum_{s=1}^{\hat S} \max_{k=1,\dots,p} \bigl(g\,c_k^{(s)}\bigr),
\]
однако в асимптотических оценках его часто опускают, поскольку \(g\) считается константой и поглощается \(O\)-большим.

\paragraph{Соглашение о вводе и выводе.}
Обычно предполагается, что входные данные не лежат целиком в памяти одного компьютера. Они считаются доступными через коммуникацию с некоторым внешним хранилищем: чтобы прочитать ввод, компьютеры получают сообщения от внешнего хранилища, а чтобы выдать результат, передают его сообщениями обратно во внешнее хранилище.

\begin{Example}{Вычисление префиксных сумм в модели BSP}{}
	
	Рассмотрим задачу вычисления префиксных сумм в модели BSP. Дана последовательность из \(n\) чисел
	\[
	x_1, x_2, \dots, x_n,
	\]
	где \(n \gg p\). Требуется вычислить последовательность
	\[
	y_1, y_2, \dots, y_n \qquad y_k = \sum_{i=1}^k x_i.
	\]
	
	Для простоты будем считать, что \(p\) делит \(n\). Тогда \(k\)-й компьютер обрабатывает блок
	\[
	x_{(k-1)n/p+1}, \dots, x_{kn/p}.
	\]
	
	\paragraph{Первый супершаг.}
	\(k\)-й компьютер считывает из внешнего хранилища свой блок длины \(n/p\).
	Вычисления не производятся, коммуникации: \(O\!\left(\frac{n}{p}\right)\).
	
	\paragraph{Второй супершаг.}
	\(k\)-й компьютер вычисляет сумму элементов своего блока и отправляет эту сумму компьютеру \(1\).
	Вычисления: \(O\!\left(\frac{n}{p}\right)\), коммуникации: \(O(p)\).
	
	\paragraph{Третий супершаг.}
	Компьютер \(1\) вычисляет префиксные суммы сумм блоков, после чего отправляет \(k\)-му компьютеру сумму всех предыдущих блоков, которую нужно прибавить к локальным префиксным суммам внутри \(k\)-го блока.
	Вычисления: \(O(p)\), коммуникации: \(O(p)\).
	
	\paragraph{Четвёртый супершаг.}
	Каждый компьютер вычисляет префиксные суммы внутри своего блока, прибавляет к ним полученную добавку и записывает результат в выходной массив.
	Вычисления: \(O\!\left(\frac{n}{p}\right)\), коммуникации: \(O\!\left(\frac{n}{p}\right)\).
	
	\paragraph{Общее время работы.}
	\begin{itemize}
		\item вычисления:
		\[
		O\!\left(p+\frac{n}{p}\right);
		\]
		
		\item коммуникации:
		\[
		O\!\left(p+\frac{n}{p}\right);
		\]
		
		\item синхронизация:
		\[
		O(1),
		\]
		так как число супершагов константно.
	\end{itemize}
	
	Если \(p^2 \ll n\), то \(p+\frac{n}{p}=O\!\left(\frac{n}{p}\right)\). Следовательно, и вычисления, и коммуникации имеют порядок \(O\!\left(\frac{n}{p}\right)\).
	
	\paragraph{Оптимальность.}
	Алгоритм асимптотически оптимален: каждое входное машинное слово нужно хотя бы один раз прочитать, поэтому требуется не менее \(\Omega\!\left(\frac{n}{p}\right)\) коммуникаций и не менее \(\Omega\!\left(\frac{n}{p}\right)\) локальных вычислений.
\end{Example}

\subsection{Модель внешней памяти}

Модель внешней памяти задаётся двумя параметрами: \(B\) и \(M\). Обычно предполагается, что
\[
B \ll M,
\]
а часто дополнительно
\[
B^2 \ll M.
\]

Алгоритм располагает неограниченной внешней памятью. За единицу времени он может прочитать или записать один блок, состоящий из \(B\) подряд идущих машинных слов внешней памяти. Кроме того, у алгоритма есть внутренняя память объёма \(M\) машинных слов, то есть \(M/B\) блоков.

Все вычисления, выполняемые над данными, уже находящимися во внутренней памяти, считаются бесплатными. Входные данные изначально расположены во внешней памяти.

\paragraph{Мотивация модели.}
На практике часто имеется процессор с относительно небольшой, но быстрой внутренней памятью и большой, но медленной внешней памятью, такой как SSD или HDD. В такой ситуации главным ограничением становится не число операций, выполняемых процессором, а объём передачи данных между внутренней и внешней памятью.

Поэтому временем работы алгоритма в этой модели считается число операций ввода-вывода, то есть число прочитанных и записанных блоков.

\subsection{Примеры задач в модели внешней памяти}

\begin{Example}{Сумма \(N\) чисел}{}
	Чтобы просуммировать \(N\) чисел, нужно хотя бы один раз прочитать весь массив. За одну операцию ввода-вывода читается блок из \(B\) элементов, поэтому время работы равно
	\[
	O\!\left(\frac{N}{B}\right).
	\]
	Эту величину часто обозначают через \(\operatorname{scan}(N)\), то есть время линейного прохода по данным.
\end{Example}

\begin{Example}{Поиск и вставка в \(B\)-дереве}{}
	Операции поиска и вставки выполняются за
	\[
	O(\log_B N),
	\]
	поскольку один узел \(B\)-дерева помещается в один блок внешней памяти, а переход на следующий уровень дерева соответствует одной операции чтения блока.
\end{Example}

\begin{Example}{Сортировка слиянием во внешней памяти}{}
	Пусть дан массив длины \(N\). Разобьём его на
	\[
	\Theta\!\left(\frac{M}{B}\right)
	\]
	подмассивов и рекурсивно отсортируем каждый из них.
	
	После этого из каждого подмассива читается по одному блоку. Так как число подмассивов равно \(\Theta(M/B)\), все эти блоки помещаются во внутреннюю память. Поэтому слияние можно выполнять, используя только данные из внутренней памяти. Когда очередной блок какого-то подмассива заканчивается, из внешней памяти дочитывается следующий блок этого подмассива за 
	\(O(1)\). В результате на всё слияние требуется
	\[
	O\!\left(\frac{N}{B}\right)
	\]
	операций ввода-вывода.
	
	Отсюда получается рекуррентная оценка
	\[
	T(N)=\Theta\!\left(\frac{M}{B}\right)\,T\!\left(\frac{N}{M/B}\right)+O\!\left(\frac{N}{B}\right).
	\]
	
	Решение этой рекуррентной формулы даёт
	\[
	T(N)=O\!\left(\frac{N}{B}\log_{M/B}\frac{N}{B}\right).
	\]
	Эквивалентно, с точностью до константы в логарифме,
	\[
	T(N)=O\!\left(\frac{N}{B}\log_{M/B}\frac{N}{M}\right).
	\]
	
	Эту величину обозначают через \(\operatorname{sort}(N).\)
\end{Example}